% GENERATED FILE. Edit canonical lessons, then run npm run generate:books.
% canonical-source-hash: fnv1a64:e8cbc5f40285266f
% canonical-lessons: FA-C134-zadan, FA-C134-tarsidan, FA-C134-chashidan, FA-C134-buyidan, FA-C134-kandan

\chapter{To hit, To be afraid, To taste, To smell, To dig}
\label{ch:fa-to-hit-to-be-afraid-to-taste-to-smell-to-dig}
\hlchaptermodality{134}

\begin{chapteropening}
\textbf{By the end of this chapter:} \emph{I can say to hit, to be afraid, to taste, to smell and to dig.}

Complete the last lesson of chapter 134: I can say to hit, to be afraid, to taste, to smell and to dig.
\end{chapteropening}

\section[\texorpdfstring{\fa{زدن} (zadan)}{زدن (zadan)}]{\fa{زدن} (zadan) — to hit}

\label{lesson:FA-C134-zadan}



\begin{quote}
Before the new one: say the Persian for to arrive, then the Persian for to weep.
\end{quote}

\subsection*{You'll want to know: \fa{زدن}}
\textbf{\fa{زدن}} — \emph{zadan} — \textquotedblleft{}to hit\textquotedblright{}.

\begin{cousinweb}
\emph{zadan} is a native Persian word, inherited rather than borrowed; it is usually traced to Middle Persian \emph{zadan}, \textquotedblleft{}to strike\textquotedblright{}. It is usually set beside Avestan \emph{jan-}, \textquotedblleft{}to strike\textquotedblright{}, and Sanskrit \emph{han-}, \textquotedblleft{}to kill\textquotedblright{}.
\end{cousinweb}

\begin{grammarlens}[title={What you've built}]
One more word for this chapter.
\end{grammarlens}

\subsection*{Your turn}
\begin{itemize}
\raggedright
  \item \emph{Say it:} \emph{zadan}
  \item \emph{Say it:} \emph{zadan}, once more
  \item \emph{Say it:} say \emph{zadan}
  \item \emph{Recall:} say the Persian for to arrive, then the Persian for to weep, then say \emph{zadan} again
\end{itemize}

\begin{tcolorbox}[breakable,colback=teal!4,colframe=teal!35!black,title={Before you move on}]
What does \fa{زدن} mean? (\textquotedblleft{}To hit\textquotedblright{}.) Say it once more.
\end{tcolorbox}
\section[\texorpdfstring{\fa{ترسیدن} (tarsidan)}{ترسیدن (tarsidan)}]{\fa{ترسیدن} (tarsidan) — to be afraid}

\label{lesson:FA-C134-tarsidan}



\begin{quote}
Before the new one: say the Persian for to weep, then the Persian for to hit.
\end{quote}

\subsection*{You'll want to know: \fa{ترسیدن}}
\textbf{\fa{ترسیدن}} — \emph{tarsidan} — \textquotedblleft{}to be afraid\textquotedblright{}.

\begin{cousinweb}
\emph{tarsidan} is a native Persian word, inherited rather than borrowed; it is usually traced to Middle Persian \emph{tarsīdan}, \textquotedblleft{}to fear\textquotedblright{}. EtymDB takes it back to the Indo-European root \emph{*tres-}, \textquotedblleft{}to tremble\textquotedblright{}, the root of Latin \emph{terrere}.
\end{cousinweb}

\begin{grammarlens}[title={What you've built}]
One more word for this chapter.
\end{grammarlens}

\subsection*{Your turn}
\begin{itemize}
\raggedright
  \item \emph{Say it:} \emph{tarsidan}
  \item \emph{Say it:} \emph{tarsidan}, once more
  \item \emph{Say it:} say \emph{tarsidan}
  \item \emph{Recall:} say the Persian for to weep, then the Persian for to hit, then say \emph{tarsidan} again
\end{itemize}

\begin{tcolorbox}[breakable,colback=teal!4,colframe=teal!35!black,title={Before you move on}]
What does \fa{ترسیدن} mean? (\textquotedblleft{}To be afraid\textquotedblright{}.) Say it once more.
\end{tcolorbox}
\section[\texorpdfstring{\fa{چشیدن} (chashidan)}{چشیدن (chashidan)}]{\fa{چشیدن} (chashidan) — to taste}

\label{lesson:FA-C134-chashidan}



\begin{quote}
Before the new one: say the Persian for to hit, then the Persian for to be afraid.
\end{quote}

\subsection*{You'll want to know: \fa{چشیدن}}
\textbf{\fa{چشیدن}} — \emph{chashidan} — \textquotedblleft{}to taste\textquotedblright{}.

\begin{cousinweb}
\emph{chashidan} is a native Persian word, inherited rather than borrowed; it is usually traced to Middle Persian \emph{čašīdan}, \textquotedblleft{}to taste\textquotedblright{}.
\end{cousinweb}

\begin{grammarlens}[title={What you've built}]
One more word for this chapter.
\end{grammarlens}

\subsection*{Your turn}
\begin{itemize}
\raggedright
  \item \emph{Say it:} \emph{chashidan}
  \item \emph{Say it:} \emph{chashidan}, once more
  \item \emph{Say it:} say \emph{chashidan}
  \item \emph{Recall:} say the Persian for to hit, then the Persian for to be afraid, then say \emph{chashidan} again
\end{itemize}

\begin{tcolorbox}[breakable,colback=teal!4,colframe=teal!35!black,title={Before you move on}]
What does \fa{چشیدن} mean? (\textquotedblleft{}To taste\textquotedblright{}.) Say it once more.
\end{tcolorbox}
\section[\texorpdfstring{\fa{بوییدن} (buyidan)}{بوییدن (buyidan)}]{\fa{بوییدن} (buyidan) — to smell}

\label{lesson:FA-C134-buyidan}



\begin{quote}
Before the new one: say the Persian for to be afraid, then the Persian for to taste.
\end{quote}

\subsection*{You'll want to know: \fa{بوییدن}}
\textbf{\fa{بوییدن}} — \emph{buyidan} — \textquotedblleft{}to smell\textquotedblright{}.

\begin{cousinweb}
\emph{buyidan} is built on \emph{bu}, \textquotedblleft{}smell\textquotedblright{}, with the verb ending \emph{-idan}. \emph{bu} is native, usually traced to Middle Persian \emph{bōy}.
\end{cousinweb}

\begin{grammarlens}[title={What you've built}]
One more word for this chapter.
\end{grammarlens}

\subsection*{Your turn}
\begin{itemize}
\raggedright
  \item \emph{Say it:} \emph{buyidan}
  \item \emph{Say it:} \emph{buyidan}, once more
  \item \emph{Say it:} say \emph{buyidan}
  \item \emph{Recall:} say the Persian for to be afraid, then the Persian for to taste, then say \emph{buyidan} again
\end{itemize}

\begin{tcolorbox}[breakable,colback=teal!4,colframe=teal!35!black,title={Before you move on}]
What does \fa{بوییدن} mean? (\textquotedblleft{}To smell\textquotedblright{}.) Say it once more.
\end{tcolorbox}
\section[\texorpdfstring{\fa{کندن} (kandan)}{کندن (kandan)}]{\fa{کندن} (kandan) — to dig}

\label{lesson:FA-C134-kandan}



\begin{quote}
Before the new one: say the Persian for to taste, then the Persian for to smell.
\end{quote}

\subsection*{You'll want to know: \fa{کندن}}
\textbf{\fa{کندن}} — \emph{kandan} — \textquotedblleft{}to dig\textquotedblright{}.

\begin{cousinweb}
\emph{kandan} is a native Persian word, inherited rather than borrowed; it is usually traced to Middle Persian \emph{kandan}, \textquotedblleft{}to dig\textquotedblright{}.
\end{cousinweb}

\begin{grammarlens}[title={What you've built}]
One more word for this chapter.
\end{grammarlens}

\subsection*{Your turn}
\begin{itemize}
\raggedright
  \item \emph{Say it:} \emph{kandan}
  \item \emph{Say it:} \emph{kandan}, once more
  \item \emph{Say it:} say \emph{kandan}
  \item \emph{Recall:} say the Persian for to taste, then the Persian for to smell, then say \emph{kandan} again
\end{itemize}

\begin{tcolorbox}[breakable,colback=teal!4,colframe=teal!35!black,title={Before you move on}]
What does \fa{کندن} mean? (\textquotedblleft{}To dig\textquotedblright{}.) Say it once more.
\end{tcolorbox}
